\documentclass[11pt]{article}
\usepackage[margin=1in]{geometry}
\usepackage{amsmath,amssymb,amsthm}
\usepackage[T1]{fontenc}
\usepackage{lmodern}
\usepackage{hyperref}
\title{A Rational Rotation Refutes the Prime-Time LIL as Stated}
\author{quantumtool}
\date{October 19, 2026}
\begin{document}
\maketitle
\section*{Counterexample under the stated quantifiers}
The statement has no restriction that $\alpha$ be irrational. Choose $\alpha=1/2$ and $f(x)=\cos(2\pi x)$ on $[0,1]$. This is a nonconstant function with mean $\int_0^1 f(x)\,dx=0$.

Every odd prime $p$ has fractional part $\{p/2\}=1/2$, so $f(\{p/2\})=-1$. The prime $2$ contributes $f(\{2/2\})=f(0)=1$. Hence for $x\ge 2$ the displayed centered sum is exactly
\[
  \sum_{p\le x} f(\{p/2\})-\pi(x)\int_0^1f
  =2-\pi(x).
\]
Since there are infinitely many primes, this is eventually negative and has no positive fluctuations at any positive scale. It therefore cannot satisfy the claimed two-sided $\Omega_{\pm}(\sqrt{x\log\log x/\log x})$ statement. If the intended conjecture assumes irrational $\alpha$, that hypothesis is absent from both displayed versions of the problem.

\section*{Scope of the formalization}
The Lean theorem verifies the arithmetic core that every prime greater than $2$ is odd, from an elementary factorization definition of primality. The fractional-part evaluation, integral of the cosine, infinitude of primes, and translation from eventual negativity to failure of $\Omega_{\pm}$ are written out here and are not formalized in Lean.
\section*{Lean formalization}
The self-contained Lean file proves that any prime at least three is congruent to one modulo two. The trigonometric and asymptotic steps are explicit in the report but are not formalized.
\end{document}
